% Source: https://artofproblemsolving.com/community/c3h1453680_statistics
% Stable identity: aops-topic-1453680-q01
\textbf{(a)} The median is a positive integer, so it is at least $1$. It can be $1$: the numbers
\[
1,1,1,1,1,203,203,204,204
\]
have sum $5+814=819=9\cdot91$, so their average is $91$, and their median (the fifth number) is $1$. The smallest possible median is $\boxed{1}$.

\textbf{(b)} \textbf{No limit.} Take the numbers $91,91,91,91,91,N,N,N,N$ with $N\ge91$. The median is $91$, and the average $\dfrac{455+4N}{9}$ can be made as large as we like by choosing $N$ large. So there is no largest possible average. (Only the bottom five numbers are restricted by the median; the four largest numbers can be arbitrarily large.)
